All three independently initialized models passed the comparison targets.
Their midpoint estimates were 26.1576, 26.1609 and 26.1589, respectively,
against the published 26.154. The absolute differences ranged from 0.0036 to
0.0069, and each joint enclosure had width about 0.0745. The source row's
published enclosure has width 0.033; the present estimates therefore achieve
close price agreement with wider sampling uncertainty.

\begin{center}
\begin{tabular}{lrrr}
\toprule
Model & Lower & Upper & Midpoint\\
\midrule
Published &26.1560&26.1520&26.1540\\
Volt A &26.1393&26.1759&26.1576\\
Volt B &26.1427&26.1790&26.1609\\
Volt C &26.1407&26.1771&26.1589\\
\bottomrule
\end{tabular}
\end{center}
The respective joint 95\% enclosures were $[26.1206,26.1950]$,
$[26.1239,26.1984]$ and $[26.1221,26.1964]$. Common assessment paths couple
the three results; the three initializations measure training variability
under that shared assessment. Every checkpoint was selected and archived
before the first assessment result was opened.

Each full CPU float64 fit completed 30,500 updates in about 770 seconds.
The three fits ran concurrently, with one CPU thread per fit. This duration
includes the complete backward training sequence and its progress checks;
independent pricing follows as a separate calculation. The terminal-only
European value is 23.0516, obtained by independent integration of the
maximum-price distribution. Early exercise adds about 3.11 in this case.

The optional fixed-translation representation passed twelve independent
controls spanning CPU float64/float32, Metal float32 and CUDA float32, with
both converted and already-fixed initial states. The original float32
representation exposed post-update discrepancies in the redundant
translations; those observations motivated the explicit alternative.


The matched float32 measurements below use the optional fixed-translation
representation and strict native-dtype products. All eighteen CPU/GPU workload
comparisons passed the numerical controls. Each entry is the median of sixty
measured updates across three restored blocks, each preceded by three warm-ups.
Update times include forward/backward calculation, Adam, synchronization and
scalar checks. Here ms means milliseconds.

\begin{center}
\setlength{\tabcolsep}{3pt}
\begin{tabular}{rrrrrr}
\toprule
& &\multicolumn{2}{c}{Apple M3 Max}&\multicolumn{2}{c}{Linux workstation}\\
$d$&Batch&CPU&Metal&CPU&CUDA\\
\midrule
5&256&0.65&7.08&1.39&5.38\\
5&1,024&2.33&9.67&4.34&3.05\\
5&8,192&16.62&38.02&38.48&11.91\\
10&256&0.73&6.79&1.50&2.89\\
10&1,024&2.44&8.79&5.45&3.14\\
10&8,192&18.57&41.94&41.97&11.19\\
50&256&1.34&8.66&2.38&2.76\\
50&1,024&4.54&10.29&9.48&3.56\\
50&8,192&34.89&56.68&77.69&16.13\\
\bottomrule
\end{tabular}
\end{center}
The Mac runs macOS 27.0.1 with 128 GiB unified memory. The Linux system uses
one CPU thread on a Threadripper PRO 7985WX and an RTX 4090 with 24 GiB device
memory, CUDA 12.5 and driver 555.42.02. Both use Rust 1.99.1 and Burn 0.22.1.
Metal fusion is enabled; CUDA uses the qualified unfused configuration.
Background host utilization was recorded throughout, including concurrent
CPU training on the Mac. These are measurements under that observed load.

CPU execution was faster than Metal throughout these training cases. CUDA
became useful at batch 1,024: the five-stock case reached a 1.42-fold update
speed ratio over its same-host CPU, increasing to 3.23-fold at batch 8,192.
The fifty-stock, batch-8,192 ratio was 4.81. A larger state and resident batch
provide enough arithmetic to amortize device dispatch and synchronization.

First-use work also matters. For five stocks and batch 8,192, the first training
update took about 629 ms on Metal and 2,039 ms on CUDA. First inference and
readback took 103 and 412 ms respectively, while warmed inference/readback took
1.62 and 0.63 ms. These first-touch figures include the existing compiler/cache
state. The full timing record separates parameter setup, upload, first use and
steady updates, making it possible to choose a device for the actual request.
